% =============
% 3.2 Use Case: Üben
% =============
\section{Use Case: Üben}

Der Kern der Anwendung ist der Übungsmodus, welcher Schüler:innen ermöglicht, gezielt
Maturafragen zu trainieren und ihr Wissen Stück für Stück auszubauen. Anders als beim
Prüfungsmodus steht hier nicht der Leistungsdruck im Vordergrund, sondern das aktive 
Lernen mit direktem Feedback. Nutzer:innen soll es ermöglicht werden, ein Fundament für 
ihr Wissen festigen zu können. 

Funktionen umfassen das Auswählen von bestimmten Themen, das Zusammenstellen 
von einem eigenen Fragenpool und das Beantworten von Fragen ohne Stress. Nutzer:innen können hierbei also genau die Themen lernen, die ihnen
Probleme bereiten, und das durch eigene Fragen oder Fragen anderer Nutzer:innen. 

Die Plattform unterstützt dabei verschiedene Fragetypen, um das Lernerlebnis flexibler zu 
gestalten. Um nicht einfach nur ein Ergebnis von ,,richti'' und,,falsch'' zu liefern, 
können Nutzer:innen auch Lösungswege einsehen, um das Verständnis nachhaltig zu fördern.

Durch diese wiederholbare und flexible Übungsumgebung eignet sich die Anwendung sowohl
für den schnellen Gebrauch als Wissenscheck als auch für die langfristige Vorbereitung
auf Tests, Schularbeiten oder die Matura.

% =============
% 3.2.1 Verwaltung von Fragenpools
% =============
\subsection{Verwaltung von Fragenpools}
\label{chap:fragenpools}
\setauthor{\secondauthor}

Die Verwaltung von Fragenpools ist eine zentrale Funktion der Lernplattform, die es
Nutzer:innen ermöglicht, ihre Lerninhalte individuell zu gestalten und gezielt zu üben. Hier spielt die 
,,Fragen browsen''-Ansicht, wie in \autoref{fig:fragen-browsen} gezeigt, eine wichtige Rolle, da diese den Nutzer:innen erlaubt, ihre individuell selektierten Fragen gesammelt zu verwalten, und somit 
genügend Raum für gezieltes Lernen bietet. 

\begin{figure}[H]
    \centering
    \includegraphics[width=0.9\textwidth]{fragenBrowsen.png}
    \caption{Ansicht der Verwaltung des Fragenpools}
    \label{fig:fragen-browsen}
\end{figure}

Das Interface der Fragenpool-Verwaltung ist so gestaltet, dass Nutzer:innen schnell und intuitiv
auf ihre gespeicherten Fragen zugreifen können. Die Ansicht ist in zwei Bereiche unterteilt: Links
befindet sich eine Seitenleiste, die als Navigation zwischen den Fragen dient, während der Hauptbereich die passenden Fragen zum gewählten Thema anzeigt.

Im Hauptbereich kann zwischen ,,bereits in Fragenpool" und "alle" gewählt werden, wobei gefiltert wird, welche Fragen angezeigt werden sollen. 
Die Fragen können jederzeit hinzugefügt oder entfernt werden, indem entsprechende Buttons wie "alle auswählen"
oder das Sternsymbol neben der einzelnen Frage verwendet werden.

Jede Frage wird in einer kompakten Kartenansicht dargestellt, die wesentliche Informationen wie die Fragestellung, dazugehörige Bilder und die Möglichkeit enthält, genau diese Frage direkt zu üben.

\subsubsection{Datenfluss in der Fragenpool-Verwaltung}

Der zentrale Teil der Fragenpool-Verwaltung ist die Interaktion mit dem Backend, um Fragen zu speichern, abzurufen und nach Fächern und Themen zu filtern.

\textbf{Navigation und Themenkatalog}

Sobald der/die Nutzer:in die Fragenpool-Verwaltung öffnet, sendet das Frontend eine Anfrage an das Backend, um den Gesamtkatalog der verfügbaren Fächer und Themen abzurufen. 
Diese Fächer und Themen werden in der Seitenleiste angezeigt, sodass der/die Nutzer:in schnell zwischen verschiedenen Themen wechseln kann, 
die dann parallel in der Hauptansicht dargestellt werden.

Die Struktur der Daten ist hierarchisch aufgebaut: Ein Subject (Fach) enthält mehrere TopicPools (Themenpools), die wiederum verschiedene Questions (Fragen) beinhalten. Diese Hierarchie wird im Frontend in verschachtelten Dropdowns dargestellt, wobei der Zustand jedes Dropdowns (geöffnet/geschlossen) in einem Objekt gespeichert wird, das die jeweiligen IDs als Schlüssel verwendet.

\textbf{Fragen abrufen und Anzeige in der Hauptansicht}

Beim Laden der Fragenpool-Verwaltung sendet das Frontend eine Anfrage an das Backend, um die Liste der Fragen abzurufen, die der/die Nutzer:in zuvor in seinen/ihren Fragenpool aufgenommen hat. 
Diese Fragen werden in einem lokalen Array gespeichert und dienen somit als Referenz für die Anzeige der bereits ausgewählten Fragen. Die Fragen, die schon im Fragenpool sind, werden visuell mittels ausgefülltem Sternsymbol hervorgehoben (siehe \autoref{fig:fragen-browsen}).

Zusätzlich kann der/die Nutzer:in zwischen zwei Modi wählen:
\begin{itemize}
    \item \textbf{Browse-Modus}: Zeigt alle verfügbaren Fragen im gesamten Katalog an.
    \item \textbf{Pool-Modus}: Zeigt nur die Fragen, die bereits im eigenen Fragenpool enthalten sind.
\end{itemize}

Das Backend antwortet mit einer strukturierten Liste von Frageobjekten, die alle notwendigen Informationen für die Anzeige im Frontend enthalten. Die Komponente verwendet eine berechnete Getter-Methode, die abhängig vom aktuellen Modus entweder die direkten Fragen (Browse-Modus) oder die aus den QuestionPoolEntry-Objekten extrahierten Fragen (Pool-Modus) zurückgibt. Diese Architektur ermöglicht es, dieselbe Komponente für beide Ansichten wiederzuverwenden.

\textbf{Hinzufügen und Entfernen von Fragen}

Wenn der/die Nutzer:in eine Frage zu seinem/ihrem persönlichen Pool hinzufügen möchte, wird durch das Klicken auf das Sternsymbol oder den "Alle auswählen"-Button eine Anfrage an das Backend gesendet. Das Backend speichert die Zuordnung zwischen Nutzer:in und Frage in der Datenbank und bestätigt die erfolgreiche Operation.

Die Implementierung dieser Funktionalität folgt dem Prinzip des Optimistic UI Updates: Sobald der/die Nutzer:in auf das Sternsymbol klickt, wird die Anzeige sofort aktualisiert, ohne auf die Antwort des Backends zu warten. Die Frage-ID wird unmittelbar zum Array der ausgewählten Fragen hinzugefügt, wodurch das Sternsymbol seinen Zustand ändert. Parallel dazu läuft die Backend-Anfrage im Hintergrund. Diese Implementierung sorgt für eine flüssige Benutzererfahrung, da keine Ladezeiten wahrgenommen werden.

Umgekehrt kann eine Frage durch erneutes Klicken auf das Sternsymbol aus dem Pool entfernt werden. Auch hier erfolgt eine sofortige Aktualisierung der UI, während die Backend-Anfrage asynchron verarbeitet wird. Das Backend löscht die entsprechende Zuordnung aus der Datenbank. Die Änderungen werden sofort im Frontend reflektiert, sodass die Nutzer:innen unmittelbares visuelles Feedback erhalten.

Für die Mehrfachauswahl existiert ein spezieller Auswahlmodus, der über einen Button aktiviert werden kann. In diesem Modus ändert sich das Verhalten der Fragenkarten: Statt die Details einer Frage anzuzeigen, werden Fragen zur Auswahlliste hinzugefügt oder entfernt. Dies ermöglicht es, mehrere Fragen gleichzeitig auszuwählen und sie dann mit einer einzigen Backend-Anfrage zum Pool hinzuzufügen (Batch-Operation). Diese Funktionalität ist besonders nützlich, wenn Nutzer:innen ihren Pool initial mit vielen Fragen befüllen möchten.

\textbf{Filterung und Suche}

Die Fragenpool-Verwaltung unterstützt verschiedene Filteroptionen, um die Übersichtlichkeit bei großen Fragenkatalogen zu gewährleisten:
\begin{itemize}
    \item \textbf{Filterung nach Fach}: Über die Seitenleiste können Nutzer:innen gezielt Fragen eines bestimmten Faches anzeigen lassen.
    \item \textbf{Filterung nach Thema}: Innerhalb eines Faches kann weiter nach spezifischen Themengebieten gefiltert werden.
    \item \textbf{Pool-Status}: Die Ansicht kann auf bereits ausgewählte Fragen beschränkt oder auf alle verfügbaren Fragen erweitert werden.
\end{itemize}

Die Filterung nach Pool-Status erfolgt clientseitig ohne erneute Backend-Anfragen. Dafür wird eine Kopie aller geladenen Fragen in einem separaten Array gespeichert. Beim Aktivieren des Filters werden nur jene Fragen angezeigt, deren IDs im Array der ausgewählten Fragen enthalten sind. Beim Deaktivieren wird die vollständige Liste wiederhergestellt. Diese Implementierung minimiert die Netzwerklast und ermöglicht sofortige Filteränderungen ohne Wartezeit.

Diese Filtermöglichkeiten erleichtern es den Nutzer:innen, gezielt die Fragen zu finden und zu üben, die für ihre aktuelle Lernphase relevant sind.

\subsubsection{Technische Umsetzung}

\textbf{Frontend-Architektur}

Die zentrale Ansicht für den Fragenpool wird von der QuestionBrowsingViewComponent gesteuert. Sie hält den aktuellen Zustand der Ansicht und greift über Services auf das Backend zu. Diese Services kümmern sich um die Kommunikation mit dem Server und liefern sauber strukturierte Daten zurück.

Der QuestionPoolService übernimmt dabei die wichtigsten Aufgaben rund um den Fragenpool: Er lädt alle Fragen eines Nutzers, zeigt Fragen zu bestimmten Themen an, fügt neue Fragen hinzu, entfernt bestehende und aktualisiert die Statistik.

\textbf{Backend-Architektur}

Wenn Fragen zum Fragenpool hinzugefügt werden, lädt das Backend zuerst den bestehenden Pool des Nutzers. Anschließend wird geprüft, ob eine Frage bereits vorhanden ist, um doppelte Einträge zu vermeiden. Nur neue Fragen werden dem Pool hinzugefügt.
Die Verbindung zwischen Fragenpool und Frage erfolgt über ein eigenes Objekt, das zusätzliche Informationen speichert. Dazu gehört, wie oft eine Frage richtig beantwortet wurde und ob sie zuletzt korrekt gelöst wurde. Diese Daten sind später wichtig für Statistiken und ein angepasstes Lernverhalten.
Alle Änderungen an der Datenbank laufen abgesichert ab: Entweder wird alles korrekt gespeichert oder bei einem Fehler vollständig zurückgesetzt. So bleiben die Daten immer konsistent.
Die Struktur der Fächer und Themen ist klar hierarchisch aufgebaut: Ein Fach enthält mehrere Themenpools, und jeder Themenpool besteht aus mehreren Fragen. Diese Struktur wird direkt in der Benutzeroberfläche dargestellt und erleichtert die Navigation.


% =============
% 3.2.2 Fragenhandling und Fragetypen
% =============
\subsection{Fragenhandling und Fragetypen}
\setauthor{\firstauthor}

Wie in \autoref{fig:question-runner-example} dargestellt befindet sich der/die Nutzer:in nun beim
Question-Runner, also der zentralen Ansicht, bei der Fragen beantwortet und ausgewertet werden können. Diese Komponente ist eine dynamische Ansicht, da sie je nach gewähltem Modus (Üben, Prüfen oder Review) verschiedene Steuerungselemente bereitstellt. 


\vspace{1cm}
\textbf{Benutzeroberfläche des Question Runners}

\begin{figure}[H]
	\centering
	\includegraphics[width=0.8\textwidth]{question-runner-mc.png}
	\caption{Beispielhafte Ansicht einer Frage im Question Runner}
	\label{fig:question-runner-example}
\end{figure}

Die Ansicht ist hierarchisch aufgebaut, um für eine schnelle Orientierung zu sorgen:
\begin{itemize}
	\item \textbf{Informationsbereich:} \\
	Am oberen Rand wird das aktuelle Fach sowie der spezifische Themenpool angezeigt.
	\item \textbf{Inhaltsbereich:} \\
	Die Fragestellung wird auffällig als Überschrift präsentiert und bei Bedarf auf der rechten Seite durch eine illustrative Grafik ergänzt.
	\item \textbf{Interaktionsbereich:} \\
	Wie nachfolgend beschrieben erscheinen hier je nach Fragetyp Auswahlfelder (Multiple-Choice) oder ein Texteingabefeld (Freitext).
	\item \textbf{Aktionsleiste:} \\
	Im unteren Bereich befinden sich die Navigationselemente. Im Übungsmodus ist dies primär der Button „Prüfen“ zur sofortigen Auswertung, während in der Prüfungssimulation (Kapitel \ref{chap:pruefung}) Navigationsbuttons zum Vor- und Zurückblättern dominieren.
\end{itemize}

\vspace{1cm}
\textbf{Unterstützung verschiedener Fragetypen}

Neben der Unterscheidung der Modi werden im Question Runner die verschiedenen Fragetypen berücksichtigt. Aktuell werden zwei Haupttypen unterstützt:

\begin{itemize}
    \item \textbf{Multiple-Choice-Fragen:} 
        Nutzer:innen wählen eine oder mehrere richtige Antworten aus einer Liste vorgegebener Auswahlmöglichkeiten aus. 
        Multiple-Choice-Fragen sind besonders geeignet, um Faktenwissen abzufragen oder schnelle Wissenskontrollen durchzuführen. 
        Das System wertet die Antwort automatisch aus und kann im Übungsmodus direkt Feedback geben. 
        Die Darstellung erfolgt als übersichtliche Auswahlfelder, die klar zwischen ausgewählt und nicht ausgewählt unterscheiden.
        (siehe \autoref{fig:question-runner-mc})
        
    \item \textbf{Freitext-Fragen:} 
        Nutzer:innen geben ihre Antwort in ein Texteingabefeld ein. 
        Dieser Fragetyp eignet sich für offene Aufgabenstellungen, bei denen die Antwort nicht auf eine vorgegebene Auswahl beschränkt ist. Für die Eingabe wird ein größeres Eingabefeld bereitgestellt, um längere Antworten komfortabel angeben zu können. Die Eingabe wird mit einer Menge an korrekten Antworten abgeglichen. Entspricht die Eingabe einer dieser Antwortmöglichkeiten, wird die Antwort als korrekt bewertet, andernfalls als falsch. 
        (siehe \autoref{fig:question-runner-freetext}).
\end{itemize}

\begin{figure}[H]
    \centering
    \begin{subfigure}[b]{0.48\textwidth}
        \centering
        \includegraphics[width=\textwidth]{question-runner-mc.png}
        \caption{Multiple-Choice-Frage}
        \label{fig:question-runner-mc}
    \end{subfigure}
    \hfill
    \begin{subfigure}[b]{0.48\textwidth}
        \centering
        \includegraphics[width=\textwidth]{question-runner-freetext.png}
        \caption{Freitext-Frage}
        \label{fig:question-runner-freetext}
    \end{subfigure}
    \caption{Question-Runner: Beispiele für verschiedene Fragetypen}
    \label{fig:question-runner}
\end{figure}

Durch diese Gestaltung unterstützt der Question Runner den/die Nutzer:in auf effektive Weise beim 
Lernen, beim Prüfen des eigenen Wissens und beim Überprüfen bereits bearbeiteter Aufgaben.  
Die modulare Darstellung der Fragen, der direkte Zugriff auf Feedback im Übungsmodus und die 
klare Trennung der Modi tragen dazu bei, dass sich der/die Nutzer:in jederzeit auf das Lernen 
konzentrieren kann, ohne von technischen Details abgelenkt zu werden.


% =============
% 3.2.3 Übungsmodus mit direktem Feedback
% =============
\subsection{Übungsmodus mit direktem Feedback}
\label{chap:uebung}
\setauthor{\firstauthor}

Beim Aufruf des Übungsmodus landet der User zunächst auf einer Übersichtsseite, auf der eine Zusammenfassung
des aktuellen Lernstands zu finden ist und die einen Einstieg in den nächsten Übungsdurchlauf ermöglicht. 
Diese Ansicht dient dazu, einen schnellen Überblick darüber zu geben, wie gut einzelne Fragen bereits 
beherrscht werden und welche Bereiche noch wiederholt werden sollten.

\begin{figure}[H]
    \centering
    \includegraphics[width=0.8\textwidth]{progress-view.png}
    \caption{Fortschrittsübersicht des Übungsmodus}
    \label{fig:progress-screen}
\end{figure}


Auf dieser Übersichtsseite werden alle Fragen der ausgewählten Themenbereiche nach ihrem bisherigen Bearbeitungsstatus 
gruppiert. Dazu zählen unter anderem:

\begin{itemize}
    \item (gelb) Fragen, die noch nie beantwortet wurden \\ (rot) Fragen, die zuletzt falsch beantwortet wurden
    \item (hellgrün) Fragen, die bereits 1× korrekt beantwortet wurden \\ zum Wiederholen gesperrt oder nicht gesperrt
    \item (grün) Fragen, die bereits 2× korrekt beantwortet wurden \\ zum Wiederholen gesperrt oder nicht gesperrt
    \item (dunkelgrün) Fragen, die ausreichend häufig richtig beantwortet wurden
\end{itemize}

Diese Einteilung hilft dem:der Nutzer:in dabei, den eigenen Lernfortschritt einzuschätzen und gezielt zu entscheiden, 
welche Fragen noch verstärktere Wiederholung benötigen. Gleichzeitig hebt sie hervor, welche Inhalte sicher beherrscht werden 
und daher im Übungsmodus seltener erscheinen. Diese Übersicht macht außerdem sichtbar, welche Fragen noch nicht wiederholt
werden sollten, da sie erst vor kurzer Zeit geübt wurden. Das sind jene Fragen, die als \textit{"fürs Wiederholen gesperrt"} 
markiert werden.

Rechts neben der Statusübersicht kann der Themenbereich gewählt werden, auf den sich sowohl die Fortschrittsübersicht als auch 
die nächste Übungssitzung beschränken soll. So kann sich der Übungsmodus auf bestimmte Fächer oder Themenbereiche eingrenzen 
lassen. Wird aber bei dieser Auswahl \textit{"Alle Fragen"} angegeben, so wird der Übungsmodus auch auf alle 
Themenbereiche im eigenen Fragenpool ausgeweitet.

\begin{figure}[H]
    \centering
    \includegraphics[width=0.8\textwidth]{progress-view-topicpool-selection.png}
    \caption{Auswahl des Themenbereichs für den nächsten Übungsdurchlauf}
    \label{fig:progress-screen-topicpool-selection}
\end{figure}

Nachdem der gewünschte Themenbereich ausgewählt wurde, kann der Übungsdurchlauf über den "Üben"-Button direkt unter der 
Auswahl des Themenbereichs gestartet werden. Das System wechselt automatisch in den Question Runner und lädt die nächste
passende Frage aus dem ausgewählten Themenbereich, je nachdem bei welchem Fragenstatus aktuell am meisten Übungsbedarf 
besteht. Auf diese Weise kann die Übungssitzung direkt an dem Punkt starten, an dem es zu diesem Zeitpunkt am wichtigsten
ist.


\vspace{1cm}
\textbf{Bearbeitung der Fragen und direktes Feedback}

Beim Start in den Übungsdurchlauf gelangt der/die Nutzer:in in den Question Runner, in dem die Fragen nacheinander 
bearbeitet werden. 
In diesem Modus liegt der Fokus auf dem Erlernen der ausgewählten Inhalte. Es ist dabei also wichtig, dass es keine 
Elemente gibt, die möglicherweise Stress oder Ablenkung verursachen, wie sie in klassischen Prüfungssituationen häufig auftreten. 
Deshalb ist diese Ansicht auch besonders schlicht und simpel gehalten. Der:Die Nutzer:in steht hier nicht unter Zeitdruck
und kann die Frage in Ruhe lesen, überdenken und schließlich beantworten. Erst nachdem eine Antwort ausgewählt 
beziehungsweise eingegeben wurde, kann diese über den \textit{"Prüfen"}-Button bestätigt werden. Anschließend wird die 
Antwort ausgewertet und das Ergebnis unmittelbar im Question Runner angezeigt.

Das direkte Feedback ist so gestaltet, dass es den Lernprozess aktiv unterstützt. Nach der Auswertung wird die Frage deshalb
visuell markiert und Antworten werden entsprechend farbig markiert. Korrekte Antworten werden mit einem positiven Grünton, 
falsche Antworten rot hervorgehoben (siehe \autoref{fig:question-answer-in-correctly}). Zusätzlich wird die richtige Lösung 
und ein möglicher Lösungsweg (siehe Kapitel \ref{chap:solutions}) angezeigt. So können Fehler sofort nachvollzogen und 
Missverständnisse direkt geklärt werden.

\begin{figure}[H]
    \centering
    \begin{subfigure}[b]{0.48\textwidth}
        \centering
        \includegraphics[width=\textwidth]{question-answered-correctly.png}
        \caption{Korrekt beantwortete Frage}
    \end{subfigure}
    \hfill
    \begin{subfigure}[b]{0.48\textwidth}
        \centering
        \includegraphics[width=\textwidth]{question-answered-incorrectly.png}
        \caption{Falsch beantwortete Frage}
    \end{subfigure}
    \caption{Question-Runner: Richtig und Falsch beantwortete Frage}
    \label{fig:question-answer-in-correctly}
\end{figure}

Durch diese Form der Rückmeldung wird das Gelernte mit der jeweils bearbeiteten Frage verknüpft. So bleibt der Lernprozess
nicht abstrakt, sondern findet direkt am Beispiel der Aufgabe statt. Zugleich ermöglicht die strukturierte Darstellung, 
dass auch falsch beantwortete Fragen nicht nur negativ und als Fehler wahrgenommen, sondern auch als positiver
Bestandteil des Lernprozesses genutzt werden.

Nachdem die Auswertung erfolgt ist, wechselt der \textit{"Prüfen"}-Button zum \textit{"Nächste Frage"}-Button. Dieser 
Button bringt den/die Nutzer:in direkt zur nächsten passenden Frage des ausgewählten Themenbereichs, wobei die Auswahl der 
nächsten Frage der in Kapitel \ref{chap:lernstrategien} beschriebenen Lernstrategie folgt. Der Übungsmodus führt die Nutzer:innen 
somit schrittweise durch die Inhalte, die für den weiteren Lernfortschritt am relevantesten sind. Dabei muss der/die Nutzer:in nicht 
selbstständig die nächste Frage auswählen, sondern kann sich einfach vom System leiten lassen.


\vspace{0.5cm}
\textbf{Wiederholungslogik und erneute Präsentation von Fragen}

Nach der Bearbeitung einer Frage wird deren Ergebnis gespeichert. Richtig beantwortete Fragen gelten zunächst als 
\textit{zum Wiederholen gesperrt} und sind erst nach einer definierten Wartezeit wieder im Übungsmodus vorzufinden. So wird
vermieden, dass dieselbe Frage unmittelbar hintereinander erneut gestellt wird, sondern erst zu einem späteren Zeitpunkt,
wenn der Lerneffekt größer ist. Falsch beantwortete Fragen bleiben im Übungsmodus, bis sie korrekt beantwortet wurden,
sodass der/die Nutzer:in direkt das Verständnis der gemachten Fehler überprüfen kann. Das soll verhindern, dass Fehler 
ein weiteres Mal gemacht werden und es zu einem endlosen Kreislauf wird, da man den Fehler nie wirklich verstehen konnte.

Fragen, die mehrfach korrekt beantwortet wurden, werden nach und nach als gefestigt eingestuft. Wurde eine Frage bereits 
dreimal hintereinander richtig beantwortet, so gilt sie als ausgelernt und wird dauerhaft aus dem aktiven Fragenpool für 
den Übungsmodus entfernt. Sie erscheint also nicht mehr im Übungsmodus, da davon ausgegangen werden kann, dass der Inhalt 
ausreichend gefestigt ist.


\vspace{0.5cm}
\textbf{Verhalten bei Sitzungsfortsetzung und Abschluss des Übungsdurchlaufs}

Das explizite Abbrechen des Übungsdurchlaufs ist nicht vorgesehen, verlässt der/die Nutzer:in den Question Runner während der 
Bearbeitung dennoch, geht der aktuelle Bearbeitungsstand verloren, die bereits beantworteten Fragen bleiben aber gespeichert.
Deshalb kehrt der/die Nutzer:in beim erneuten Aufrufen des Übungsmodus also auch automatisch zur zueletzt offenen Frage zurück.

Sind für den ausgewählten Themenbereich vorübergehend keine weiteren Fragen mehr verfügbar, weil alle Fragen entweder
ausgelernt oder zur Wiederholung gesperrt sind, so wird der \textit{"Prüfen}-Button zum \textit{"Abschließen"}-Button. Durch
diesen gelangt man zurück zur vorherigen Ansicht, also der Fortschrittsübersicht.

\begin{figure}[H]
    \centering
    \includegraphics[width=0.8\textwidth]{check-button-to-end.png}
    \caption{"Prüfen"-Button wurde zum "Abschließen"-Button}
\end{figure}


\subsection{Lösungswege und Nutzerinteraktion}
\label{chap:solutions}
\setauthor{\secondauthor}
\textit{Einsehen von Lösungswegen, Feedbackmechanismen}

